% Reproduce los casos difíciles de las notas de Álgebra Lineal (Uniandes):
% figura TikZ con etiquetas, figura al costado del texto (wrapfig), \vv{PQ} (esvect),
% vectores en línea, producto interno con ⟨ ⟩, matriz en bloque, sistema alineado y ≠.
\documentclass{article}
\usepackage[utf8]{inputenc}
\usepackage[spanish]{babel}
\usepackage{amsmath,amssymb,tikz,wrapfig,esvect}
\usetikzlibrary{babel} % babel en español activa < y >, que rompen las flechas de TikZ
\begin{document}
\section{Vectores en $\mathbb{R}^2$}
\begin{wrapfigure}{r}{0.4\textwidth}
\begin{tikzpicture}
\draw[help lines] (0,0) grid (4,4);
\draw[->] (0,0) -- (4.3,0) node[right] {$x$};
\draw[->] (0,0) -- (0,4.3) node[above] {$y$};
\draw[->,thick,blue] (1,1) -- (3,3) node[right] {$Q$};
\node at (1,0.7) {$P$};
\end{tikzpicture}
\caption{El vector $\vv{PQ}$ y una de sus representaciones.}
\end{wrapfigure}
Dados dos puntos $P, Q$ en el plano $xy$, escribimos $\vv{PQ}$ para el vector representado por la flecha que empieza en $P$ y termina en $Q$. Por ejemplo, sean $P(2,1)$ y $Q(4,4)$ puntos en el plano. Denotamos el vector por $\vv{OP}$ o $\begin{pmatrix} p_1 \\ p_2 \end{pmatrix}$ o a veces por $(p_1,p_2)^t$ para ahorrar espacio. Un vector muy especial es el vector cero $\vec 0$. Esta frase sigue al lado de la figura para comprobar que el texto que la rodea no se pierde en la conversión.

\begin{figure}[h]
\centering
\begin{tikzpicture}
\draw[help lines] (-1,-1) grid (6,6);
\draw[->] (-1,0) -- (6.5,0) node[right] {$g$};
\draw[->] (0,-1) -- (0,6.5) node[above] {$r$};
\draw[violet] (-1,6) -- (6,-1);
\node[violet] at (9,3) {$L_1 : y = 5 - x$};
\end{tikzpicture}
\caption{Gráfica de la recta $L_1$.}
\end{figure}

Ejemplo: $\left\langle \begin{pmatrix} 2 \\ 3 \end{pmatrix}, \begin{pmatrix} -1 \\ 5 \end{pmatrix} \right\rangle = 2 \cdot (-1) + 3 \cdot 5 = 13$. Si $\beta \neq 0$, entonces la recta tiene pendiente $-\frac{\alpha}{\beta}$.
\[ A = \begin{pmatrix} a_{11} & a_{12} & \cdots & a_{1n} \\ \vdots & & & \vdots \\ a_{m1} & a_{m2} & \cdots & a_{mn} \end{pmatrix} \]
\begin{align*}
x + 2y &= 11, \\
3x + 4y &= 27.
\end{align*}
\end{document}
